\section{Model}\label{sec:model}

We build on the tractable heterogeneous-agent New Keynesian (THANK) framework
of \citet{Bilbiie2021} and \citet{Bilbiie2025}. The economy is populated by a
continuum of households of measure one, a continuum of monopolistically
competitive firms, and a government. Time is discrete and indexed by
$t=0,1,2,\dots$.

\paragraph{Relationship to TANK, THANK, and full HANK.}
It is useful to situate our framework relative to three strands of the
heterogeneous-agent literature before proceeding, as they differ in important
ways that bear on our results.

The Two-Agent New Keynesian (TANK) model, developed by \citet{Gali2007} and
\citet{Bilbiie2021}, populates the economy with two permanent household types:
unconstrained savers who smooth consumption according to the permanent-income
hypothesis, and \HtM{} households who consume all their income each period.
Household type is fixed and deterministic -- a household is either always a
saver or always \HtM. TANK is highly tractable and delivers clean analytical
results, but it abstracts from three empirically important features:
idiosyncratic income risk, precautionary saving, and the possibility of
transitioning between constrained and unconstrained states over the lifecycle.

The THANK model of \citet{Bilbiie2025} extends TANK by introducing stochastic
switching between states. In each period, a saver faces a positive probability
of becoming \HtM{} next period, and vice versa. This seemingly small
modification has three important consequences. First, it introduces genuine
idiosyncratic risk -- even an unconstrained saver today may be constrained
tomorrow. Second, because this risk is uninsurable under incomplete markets, it
generates a precautionary saving motive that makes savers accumulate more wealth
than the pure permanent-income hypothesis would predict, producing more
realistic MPCs. Third, it delivers empirically realistic intertemporal marginal
propensities to consume (iMPCs) that match the micro evidence on how households
spread spending responses over time \citep{Fagereng2021}. THANK thus serves as
an analytically tractable bridge between the simple TANK intuition and the full
quantitative HANK model.

The full HANK model -- as in \citet{KaplanMollViolante2018} -- replaces the two
discrete household types with a continuous distribution of households that
differ in wealth and idiosyncratic productivity. The wealth distribution is
endogenously determined by household saving decisions under incomplete markets.
This framework is quantitatively rich and matches the empirical wealth
distribution precisely, but it requires numerical solution and sacrifices the
analytical tractability needed to derive closed-form multiplier expressions.

Our strategy exploits the complementary strengths of these frameworks. We use
THANK to derive the analytical propositions in Section~\ref{sec:analytical},
where the stochastic switching structure allows us to capture idiosyncratic
risk and precautionary saving while retaining closed-form solutions. We then
take the model to a full HANK calibration in Section~\ref{sec:quant}, where the
continuous wealth distribution disciplines the quantitative magnitudes and
allows us to study how the complementarity channel varies across the wealth
distribution -- a dimension that THANK, by construction, cannot speak to.

\subsection{Households}\label{subsec:households}
There are two types of households: unconstrained savers (type $S$) and \HtM{}
households (type $H$). A fraction $\lambda\in(0,1)$ of households are \HtM{} in
any given period, and the remaining fraction $1-\lambda$ are savers. Households
switch between states stochastically: a fraction $s\in(0,1)$ of savers remain
savers next period, and a fraction $h\in(0,1)$ of \HtM{} households remain
\HtM{} next period. In steady state, the fraction of \HtM{} households is
$\lambda=(1-s)/[(1-s)+(1-h)]$, determined by the switching probabilities. This
switching structure introduces idiosyncratic risk and precautionary saving
motives absent in TANK, key features of quantitative HANK models
\citep{Bilbiie2025}.

\subsubsection*{Preferences -- general non-separable specification.}
Each household type $j\in\{S,H\}$ has preferences over private consumption
$c^{j}_{t}$, government spending $g_{t}$, and labor $n^{j}_{t}$ represented by a
general twice-continuously-differentiable utility function
$U^{j}(c^{j}_{t},g_{t},n^{j}_{t})$. We impose: $U^{j}_{c}>0$, $U^{j}_{cc}<0$
(positive, diminishing marginal utility of consumption); $U^{j}_{n}<0$,
$U^{j}_{nn}<0$ (disutility of labor); and $U^{j}_{cn}=0$ (separability between
consumption and labor, standard in NK models). The key restriction governing
the fiscal-transmission channel is
\begin{equation}\label{eq:Ucg-signs}
  U_{cg}^{j}\ \gtreqless\ 0
  \begin{cases}
    >0 & \text{Edgeworth complements: }g\text{ raises MU of }c\\
    =0 & \text{separable: \bilbiie{} baseline}\\
    <0 & \text{Edgeworth substitutes: }g\text{ lowers MU of }c
  \end{cases}
\end{equation}
We allow this to differ across types, capturing the empirically motivated idea
that wealthier unconstrained households substitute public for private services
($U_{cg}^{S}<0$), while constrained households rely on public provision to
complement their private spending ($U_{cg}^{H}>0$).

\paragraph{Steady-state utility curvature elasticities.}
All effects of non-separability on equilibrium allocations are summarized by two
type-specific steady-state elasticities. Let overbars denote steady-state
values. Define
\begin{equation}\label{eq:elasticities}
  \epsilon^{j}_{cc}\equiv-\frac{U^{j}_{cc}\bar c}{\bar U^{j}_{c}}>0,
  \qquad
  \epsilon^{j}_{cg}\equiv-\frac{U^{j}_{cg}\bar g}{\bar U^{j}_{c}}.
\end{equation}
$\epsilon^{j}_{cc}$ is the consumption Euler elasticity (equals $\sigma$ in
separable CRRA). $\epsilon^{j}_{cg}$ is the cross-elasticity of marginal utility
with respect to $g$: negative under complementarity ($U_{cg}^{j}>0$), positive
under substitutability ($U_{cg}^{j}<0$), zero under separability. These two
numbers -- evaluated at steady state for each type -- are the only objects from
the utility function that enter the equilibrium conditions.

\begin{assumption}[General non-separable preferences]\label{asm:general}
The utility function $U^{j}(c,g,n)$ satisfies the sign conditions above. The
heterogeneous-internalization case has $\epsilon^{S}_{cg}>0$ (savers
substitute: government spending reduces their marginal utility of private
consumption) and $\epsilon^{H}_{cg}<0$ (\HtM{} complement: government spending
raises their marginal utility of private consumption). The symmetric cases
$\epsilon^{S}_{cg}=\epsilon^{H}_{cg}$ and the separable case
$\epsilon^{j}_{cg}=0$ are nested as restrictions.
\end{assumption}

\paragraph{Linearized optimality conditions.}
Log-linearizing the intratemporal FOC $-U^{j}_{n}=w_{t}U^{j}_{c}$ and the saver
Euler equation $U^{S}_{c}(c^{S}_{t},g_{t})=\beta R_{t}\mathbb{E}_{t}
U^{S}_{c}(c^{S}_{t+1},g_{t+1})$ (see Appendix~\ref{app:general-nonsep} for full
derivations):
\begin{align}
  \phi\,\hat n^{j}_{t} &= \hat w_{t} - \epsilon^{j}_{cc}\hat c^{j}_{t}
    - \epsilon^{j}_{cg}\hat g_{t},\qquad j\in\{S,H\}, \label{eq:intratemp}\\
  \hat c^{S}_{t} &= \mathbb{E}_{t}\hat c^{S}_{t+1}
    -\frac{1}{\epsilon^{S}_{cc}}\hat r_{t}
    -\frac{\epsilon^{S}_{cg}}{\epsilon^{S}_{cc}}(\hat g_{t}-\mathbb{E}_{t}\hat g_{t+1}).
    \label{eq:euler}
\end{align}
These are the only two equations that differ from \bilbiie{}. Budget constraints,
market clearing, production, and profit equations are unchanged.

\begin{remark}[Two channels of non-separability]\label{rem:two-channels}
The term $-\epsilon^{j}_{cg}\hat g_{t}$ in \eqref{eq:intratemp} is the
\emph{labor-supply channel}: under $\epsilon^{j}_{cg}<0$ (complements), a rise
in $g$ raises the marginal utility of consumption, making wages more valuable
and inducing more labor supply. The term
$-(\epsilon^{S}_{cg}/\epsilon^{S}_{cc})(\hat g_{t}-\mathbb{E}_{t}\hat g_{t+1})$
in \eqref{eq:euler} is the \emph{saver direct-demand channel}: under
$\epsilon^{S}_{cg}<0$, a positive spending shock (which is expected to decay
under AR(1): $\hat g_{t}-\mathbb{E}_{t}\hat g_{t+1}=(1-\rho_{g})\hat g_{t}>0$)
directly stimulates saver consumption today. Both channels are absent in the
separable case $\epsilon^{j}_{cg}=0$.
\end{remark}

\paragraph{Savers.}
Savers maximize
$\mathbb{E}_{0}\sum_{t=0}^{\infty}\beta^{t}U^{S}(c^{S}_{t},g_{t},n^{S}_{t})$
subject to
\begin{equation}\label{eq:saver-budget}
  c^{S}_{t}+b_{t}\le w_{t}n^{S}_{t}+R_{t-1}b_{t-1}+T^{S}_{t}+\tau_{D}D^{S}_{t}.
\end{equation}
The linearized optimality conditions are \eqref{eq:intratemp} and
\eqref{eq:euler} for $j=S$.

\paragraph{Hand-to-mouth households.}
\HtM{} households consume their disposable income each period:
\begin{equation}\label{eq:htm-budget}
  c^{H}_{t}=w_{t}n^{H}_{t}+T^{H}_{t}-t^{H}_{t},
\end{equation}
where $T^{H}_{t}$ includes profit redistribution and $t^{H}_{t}=(\mu/\lambda)g_{t}$
is their tax burden under the balanced budget. The linearized intratemporal
condition is \eqref{eq:intratemp} for $j=H$.

\begin{assumption}[Linear effective consumption -- parametric special case]
\label{asm:linear}
The utility function takes the CRRA-linear-aggregator form:
\begin{equation}\label{eq:linear-util}
  U^{j}(\ctilde^{j}_{t},n^{j}_{t})=\frac{(\ctilde^{j}_{t})^{1-\sigma}}{1-\sigma}
  -\frac{(n^{j}_{t})^{1+\phi}}{1+\phi},\qquad
  \ctilde^{j}_{t}\equiv c^{j}_{t}+\theta^{j}g_{t},
\end{equation}
with $\sigma>0$, $\phi>0$, $\theta^{j}\in\mathbb{R}$. Under the zero-profit
normalization $\bar c^{H}=\bar c^{S}=\bar c$, this delivers closed-form
expressions for the elasticities:
\begin{equation}\label{eq:linear-elasticities}
  \epsilon^{j}_{cc}=\sigma\alpha^{j},\qquad
  \epsilon^{j}_{cg}=\sigma\Gamma^{j},
\end{equation}
where
$\alpha^{j}\equiv\bar c/(\bar c+\theta^{j}\bar g)$ and
$\Gamma^{j}\equiv\theta^{j}\bar g/(\bar c+\theta^{j}\bar g)=1-\alpha^{j}$.
Complementarity corresponds to $\theta^{j}<0$ ($\epsilon^{j}_{cg}<0$),
substitutability to $\theta^{j}>0$ ($\epsilon^{j}_{cg}>0$), and separability to
$\theta^{j}=0$ ($\epsilon^{j}_{cg}=0$). The ratio
$\epsilon^{j}_{cg}/\epsilon^{j}_{cc}=\Gamma^{j}/\alpha^{j}\equiv\psi^{j}
=\theta^{j}\bar g/\bar c$ is the complementarity--wealth ratio.
\end{assumption}

All results in Section~\ref{sec:analytical} are stated first under
Assumption~\ref{asm:general} in terms of $(\epsilon^{j}_{cc},\epsilon^{j}_{cg})$,
then specialized to Assumption~\ref{asm:linear} using
\eqref{eq:linear-elasticities}. The four cases nested by
Assumption~\ref{asm:linear} are: (i) \emph{symmetric complementarity}:
$\theta^{S}=\theta^{H}=\theta<0$; (ii) \emph{symmetric substitutability}:
$\theta^{S}=\theta^{H}=\theta>0$; (iii) \emph{separability}:
$\theta^{S}=\theta^{H}=0$ \citep{Bilbiie2025}; (iv) \emph{heterogeneous
internalization}: $\theta^{S}>0$, $\theta^{H}<0$.

\paragraph{Aggregate consumption.}
Aggregate consumption is a weighted average:
\begin{equation}\label{eq:agg-consumption}
  C_{t}=\lambda c^{H}_{t}+(1-\lambda)c^{S}_{t}.
\end{equation}

\subsection{Firms}\label{subsec:firms}
A representative final-good firm aggregates a continuum of differentiated
intermediate goods using a CES technology with elasticity of substitution
$\varepsilon>1$, yielding demand $y_{t}(j)=(p_{t}(j)/P_{t})^{-\varepsilon}Y_{t}$
and aggregate price index
$P_{t}=\bigl[\int_{0}^{1}p_{t}(j)^{1-\varepsilon}dj\bigr]^{1/(1-\varepsilon)}$.

Intermediate-goods firm $j$ produces using only labor: $y_{t}(j)=n_{t}(j)$
(linear technology, so marginal cost equals the real wage: $mc_{t}=w_{t}$).
Rather than Calvo staggering, we adopt \citet{Rotemberg1982} quadratic
price-adjustment costs. Each firm pays a real cost
$\tfrac{\psi_{p}}{2}\bigl[\tfrac{p_{t}(j)}{p_{t-1}(j)}-1\bigr]^{2}Y_{t}$ to
change its price. The firm's problem delivers the Rotemberg Phillips curve:
\begin{equation}\label{eq:nkpc}
  \pi_{t}=\beta\mathbb{E}_{t}\pi_{t+1}+\kappa\hat{mc}_{t}
        =\beta\mathbb{E}_{t}\pi_{t+1}+\kappa\hat w_{t},
\end{equation}
where $\kappa=(\varepsilon-1)/\psi_{p}$ and $\hat{mc}_{t}=\hat w_{t}$ from
linear production. The government implements an employment subsidy
$\tau_{s}=1/(\varepsilon-1)$ to eliminate the steady-state markup distortion,
so $\bar{mc}=1$, $\bar D=0$, and $\bar c^{H}=\bar c^{S}=\bar c$.

\begin{remark}[Rotemberg vs.\ Calvo]
Rotemberg pricing is isomorphic to Calvo at first order -- both deliver the same
log-linearized NKPC -- but is analytically cleaner because it avoids tracking
the distribution of price setters and delivers the cost directly in the firm's
static problem.
\end{remark}

\subsection{Government}\label{subsec:government}
The government purchases $G_{t}$ units of the consumption good each period,
financed by lump-sum taxes $T_{t}$ and one-period nominal bonds. The government
budget constraint is
\begin{equation}\label{eq:govt-bc}
  G_{t}+R_{t-1}B_{t-1}=T_{t}+B_{t}.
\end{equation}
For the analytical results, we assume a balanced budget $T_{t}=G_{t}$ to isolate
the pure spending multiplier from redistribution effects. Government spending
follows an AR(1):
\begin{equation}\label{eq:g-ar1}
  \hat g_{t}=\rho_{g}\hat g_{t-1}+\varepsilon^{g}_{t},\qquad \rho_{g}\in(0,1).
\end{equation}
The central bank sets the nominal interest rate according to a Taylor rule:
\begin{equation}\label{eq:taylor}
  \hat\iota_{t}=\phi_{\pi}\pi_{t}+\phi_{y}\hat y_{t}.
\end{equation}
For the analytical multiplier derivation, we consider the case of a fixed real
interest rate ($\hat r_{t}=0$), corresponding either to the zero lower bound or
an interest-rate peg. This isolates the pure demand-amplification mechanism.

\subsection{Market clearing}\label{subsec:mc}
Goods market clearing requires
\begin{equation}\label{eq:mc-goods}
  \hat y_{t}=\hat c_{t}+\hat g_{t}.
\end{equation}
Bond market clearing requires $B_{t}=0$ in equilibrium (zero net supply).
